\section{The $p$-level tower}

The general correction is generated by powers of the operator, with the sign of Remark~\ref{rem:sign}:
\begin{equation}\label{eq:tower} C^{(\le p)}_j \;=\; +\sum_{\ell=1}^{p} \frac{\dto^\ell}{(\ell+1)!}
\bigl(D\mathcal{D}_2 + \alpha\bigr)^{\ell} u^n_j,
\end{equation} where $\mathcal{D}_2$ is a spatial operator whose accuracy must be raised as $\ell$ grows. Spatial discretization error breaks the binomial symmetry of the time expansion, giving the stencil-scaling bound:

\begin{proposition}[Stencil bound; sketch]\label{prop:stencil}
Global order $2p+2$ requires corrections evaluated on at least a
$(2p+3)$-point stencil: $p=1$ needs the 3-point main stencil plus a 5-point
evaluation inside $C^{(1)}$ when full fourth-order spatial accuracy is demanded
of the correction itself; $p=2$ needs a 7-point stencil to represent the
uncancelled $u_{xxxxxx}$ residual.
\end{proposition}
\noindent\emph{Status:} \tagO for $p\ge2$ numerics (no harness yet; the bound's
derivation is a consistency argument in \texttt{OTS-NIDC.md} \S6). The $p=1$
row is \tagV via Table~\ref{tab:scalar}.

\section{Notation dictionary and reproducibility}\label{app:notation}

\subsection{Global notation dictionary}\label{app:notation-dict}

This was useful and totally AI generated! Consider giving this a look if any notation is not clear.

\paragraph{Steps.} The optimal (Chu) step is
$\dto=\dx^2/(6D)$ for diffusivity $D$; for species $i$,
$\dto^{(i)}=\dx^2/(6D_i)$, also written $\tau_i$ (and simply $\tau$ when the
scalar problem is meant). A forced or actual step is $h$; the shared two-species
step of \S\ref{sec:regime1} is $h=\dto^{(2)}$. The macro step of
\S\ref{sec:regime2} is $H=\dto^{(2)}$, decomposed as $n_r$ full steps of
$\dto^{(1)}$ plus a remainder $\rho=\theta\,\dto^{(1)}$. The local truncation
error is always $\mathcal{T}$ (never $\tau$), as in \eqref{eq:lte}.

\paragraph{Dimensionless groups.} The master variables are the \emph{step
ratio} $r=h/\tau=6D\dt/\dx^2$ (legacy diffusion number $\mu=r/6$, used only in
classical comparisons) and the \emph{step-independent reaction strength}
$\eta=\alpha\dx^2/D$. The per-step reaction exponent is $a=\alpha h=r\eta/6$
(so $a=\eta/6$ at the optimum). The mismatch parameter is
$\varepsilon=h-\tau=(r-1)\tau$; for the remainder step of Method A,
$\varepsilon_\rho=\rho-\dto^{(1)}=(\theta-1)\dto^{(1)}$. The subcycle count is
$n_r=\lfloor r\rfloor$ and the remainder fraction is $\theta=r-n_r$.

\paragraph{Amplification factors.} With $s=\sin^2(\kappa\dx/2)$:
$G_{\mathrm{FE}}$ is the factor of the full (diffusion + reaction) uncorrected
forward-Euler scheme, $G_{\mathrm{FE}}=1-\frac{2r}{3}s+\frac{r\eta}{6}$
(\S\ref{sec:method}); $G_{\mathrm{NIDC}}$ is the factor of the corrected scheme
with the $C^{(1)}$ coefficients derived at the optimum,
$G_{\mathrm{NIDC}}=G_{\mathrm{FE}}-\frac{\eta}{9}s+\frac{\eta^2}{72}$
(\S\ref{sec:method}); $G_{\mathrm{mis}}$ is the \emph{diffusion-only} factor of
the mismatch-corrected diffusion step,
$G_{\mathrm{mis}}=1-\frac{2r}{3}s+\frac{2r(r-1)}{9}s^2$ (\S\ref{sec:regime1});
$G_{\mathrm{NIDC},h}$ is the full factor of the forced-step corrected scheme,
$G_{\mathrm{NIDC},h}=G_{\mathrm{FE}}+\frac{2r(r-1)}{9}s^2-\frac{2ra}{3}s+\frac{a^2}{2}$
(\S\ref{sec:regime1}, Remark~\ref{rem:gfac}). $E(s)$ is the exact semi-discrete
one-step factor $e^{-\frac{2r}{3}s+a}$; $E_q$ (\S\ref{sec:regime2}) is
unrelated --- the degree-$q$ time extrapolant.

\paragraph{Bridge identities.} $r=6\mu$; $a=\mu\eta=r\eta/6$; $\eta=6a$ at
$r=1$; $\mu=\frac16$ (equivalently $r=1$) at the optimal step. At $r=1$ the
factors agree: $G_{\mathrm{NIDC},h}=G_{\mathrm{NIDC}}$ and the diffusion part
of $G_{\mathrm{NIDC}}$ equals $G_{\mathrm{mis}}$. Off the optimum,
$G_{\mathrm{NIDC}}$ (frozen $C^{(1)}$ coefficients) and $G_{\mathrm{NIDC},h}$
($h$-rescaled corrections plus mismatch) are different conventions; the paper's
stability statements use the former in \S\ref{sec:method} and the latter in
\S\ref{sec:regime1}, and both satisfy $|G|\le e^{a}$ on their respective
admitted ranges.

\paragraph{Indices and spaces.} Nodes are $j$ (subscript); species are $i$,
with reaction matrix $R_{i\ell}$ (the scalar-analysis reaction strength
$\alpha$ is the diagonal eigen-reaction of the moment hierarchy); spatial
wavenumber is $\kappa$; time level is $n$, the final time level is $n_f$, the
subcycle count is $n_r$, and macro-step counter is $k$ ($t^k$). Grid size is
$N_x$. The spatial domain of the companion \secStochX{} is $\Omega$ (not $D$,
which is reserved for diffusivity); its Hilbert space is $H=L^2(\Omega)$; the
Mercer eigenvalues of the covariance are $\nu_p$ (not $\mu_p$, avoiding the
diffusion-number collision). The noise channel index is $k$, the covariance
mode index $p$, and the Fourier mode index $m$ --- $k$ and $m$ are local to
the companion \secStochX{} and never appear there as steps or nodes.

\subsection{Reproducibility}\label{app:repro}

The convergence harness of Table~\ref{tab:scalar} is
\texttt{code/implementation/convergence\_test.py} in the repository, with
results logged in \texttt{code/implementation/results.md} (2026-07). The
stability claims of \S\ref{sec:method} and \S\ref{sec:regime1} were
cross-checked by independent re-derivation of every amplification factor
symbolically from the update formulas and by dense numerical sweeps over
$(r,\eta,s)$ --- zero violations of the stated bounds on the admitted ranges;
the check scripts live outside the repository tree and should be re-homed
under \texttt{code/verification/} in the next release.

\section{Open items}\label{sec:open}

\begin{itemize}
\item \textbf{Coupled discrete stability}\label{open:stab}: the
off-diagonal reaction coupling and its interpolated subcycled variant have no
independent stability proof; only the scalar building blocks of
\S\ref{sec:method} and \S\ref{sec:regime1} are established.
\item \textbf{Cross-term correction with $\partial_t u^{(2)}$} inside the
sub-steps of \S\ref{sec:regime2} (supply via finite difference or Hermite
derivative; check it does not lower the order).
\item \textbf{More than two species}: define the ordering of the ratios and
whether to nest the subcycles.
\item \textbf{$p\ge 2$}: interpolation degree $q \ge p$ and extrapolation
degree $q \ge p+1$ are predicted by the same counting, but untested.
\end{itemize}
